
\documentclass{article}
\usepackage{colortbl}
\usepackage{makecell}
\usepackage{multirow}
\usepackage{supertabular}

\begin{document}

\newcounter{utterance}

\centering \large Interaction Transcript for game `hot\_air\_balloon', experiment `air\_balloon\_survival\_el\_complexity\_hard', episode 3 with deepseek{-}v4{-}pro.
\vspace{24pt}

{ \footnotesize  \setcounter{utterance}{1}
\setlength{\tabcolsep}{0pt}
\begin{supertabular}{c@{$\;$}|p{.15\linewidth}@{}p{.15\linewidth}p{.15\linewidth}p{.15\linewidth}p{.15\linewidth}p{.15\linewidth}}
    \# & $\;$A & \multicolumn{4}{c}{Game Master} & $\;\:$B\\
    \hline

    \theutterance \stepcounter{utterance}  
    & & \multicolumn{4}{p{0.6\linewidth}}{
        \cellcolor[rgb]{0.9,0.9,0.9}{
            \makecell[{{p{\linewidth}}}]{
                \texttt{\tiny{[P1$\langle$GM]}}
                \texttt{Συμμετέχετε σε ένα συνεργατικό παιχνίδι διαπραγματεύσεων.} \\
\\ 
\texttt{Μαζί με έναν άλλο συμμετέχοντα, πρέπει να συμφωνήσετε σε ένα ενιαίο σύνολο αντικειμένων που θα κρατηθούν. Ο καθένας από εσάς έχει τη δική του άποψη για το πόσο σημαντικό είναι το κάθε αντικείμενο για εσάς (σημασία). Δεν γνωρίζετε πώς εκτιμά τα αντικείμενα ο άλλος συμμετέχων. Επιπλέον, σας δίνεται η προσπάθεια που απαιτεί κάθε αντικείμενο.} \\
\texttt{Μπορείτε να συμφωνήσετε σε ένα σύνολο μόνο εάν η συνολική προσπάθεια των επιλεγμένων αντικειμένων δεν υπερβαίνει ένα κοινό όριο:} \\
\\ 
\texttt{LIMIT = 5638} \\
\\ 
\texttt{Εδώ είναι οι τιμές της προσπάθειας των επιμέρους αντικειμένων:} \\
\\ 
\texttt{Item effort = \{"C08": 110, "C55": 727, "A43": 527, "A67": 855, "B42": 942, "A62": 617, "B19": 301, "C43": 130, "C73": 916, "C10": 212, "B36": 146, "A18": 559, "C61": 933, "C67": 740, "C42": 33, "C70": 799, "B40": 324, "C39": 841, "A21": 925, "C16": 639\}} \\
\\ 
\texttt{Εδώ είναι η προσωπική σας άποψη για τη σημασία κάθε αντικειμένου:} \\
\\ 
\texttt{Item importance values = \{"C08": 204, "C55": 821, "A43": 621, "A67": 949, "B42": 1036, "A62": 711, "B19": 395, "C43": 224, "C73": 1010, "C10": 306, "B36": 240, "A18": 653, "C61": 1027, "C67": 834, "C42": 127, "C70": 893, "B40": 418, "C39": 935, "A21": 1019, "C16": 733\}} \\
\\ 
\texttt{Στόχος:} \\
\\ 
\texttt{Στόχος σας είναι να διαπραγματευτείτε ένα κοινό σύνολο αντικειμένων που να σας ωφελεί όσο το δυνατόν περισσότερο (δηλαδή να μεγιστοποιεί τη συνολική σημασία για ΕΣΑΣ), παραμένοντας εντός του LIMIT. Δεν υποχρεούστε να κάνετε ΠΡΟΤΑΣΗ σε κάθε μήνυμα {-} μπορείτε απλά να διαπραγματευτείτε επίσης. Όλες οι τακτικές επιτρέπονται!} \\
\\ 
\texttt{Πρωτόκολλο αλληλεπίδρασης:} \\
\\ 
\texttt{Μπορείτε να χρησιμοποιήσετε μόνο τις ακόλουθες δομημένες μορφές σε ένα μήνυμα:} \\
\\ 
\texttt{ΠΡΟΤΑΣΗ: \{'A', 'B', 'C', …\}} \\
\texttt{Προτείνετε να κρατήσετε ακριβώς αυτά τα αντικείμενα.} \\
\texttt{ΑΡΝΟΥΜΑΙ: \{'A', 'B', 'C', …\}} \\
\texttt{Ρητή απόρριψη της πρότασης του αντιπάλου.} \\
\texttt{ΕΠΙΧΕΙΡΗΜΑ: \{'...'\}} \\
\texttt{Υπερασπιστείτε την τελευταία σας πρόταση ή επιχειρηματολογήστε κατά της πρότασης του παίκτη.} \\
\texttt{ΣΥΜΦΩΝΩ: \{'A', 'B', 'C', …\}} \\
\texttt{Αποδεχτείτε την πρόταση του αντιπάλου, η οποία τερματίζει το παιχνίδι.} \\
\\ 
\\ 
\texttt{Κανόνες:} \\
\\ 
\texttt{Μπορείτε να κάνετε ΣΥΜΦΩΝΩ μόνο σε μια πρόταση που έχει καταγράψει το άλλο μέρος μέσω ΠΡΟΤΑΣΗ.} \\
\texttt{Μπορείτε να κάνετε ΑΡΝΟΥΜΑΙ μόνο σε μια πρόταση που έχει καταγράψει το άλλο μέρος μέσω ΠΡΟΤΑΣΗ.} \\
\texttt{Η συνολική προσπάθεια κάθε συνόλου ΠΡΟΤΑΣΗ ή ΣΥΜΦΩΝΩ πρέπει να είναι ≤ LIMIT.} \\
\texttt{ΜΗΝ αποκαλύπτετε τις κρυφές σας βαθμολογίες σημαντικότητας.} \\
\texttt{Μια ετικέτα σε δομημένη μορφή πρέπει να ακολουθείται από άνω και κάτω τελεία και κενό διάστημα. Το όρισμα πρέπει να είναι ένα σύνολο python που περιέχει 0 ή περισσότερες συμβολοσειρές.} \\
\texttt{Έτσι, πρέπει να έχει τη μορφή "TAG: \{...\}"} \\
\texttt{Ακολουθήστε αυστηρά το πρωτόκολλο αλληλεπίδρασης και ΜΗΝ γράφετε τίποτα πέρα από τη δοσμένη δομή.} \\
\texttt{Το παιχνίδι τελειώνει όταν η μία πλευρά δώσει ένα ΣΥΜΦΩΝΩ σε ένα ΠΡΟΤΑΣΗ που έκανε ο άλλος παίκτης.} \\
\texttt{Το περιεχόμενο της απάντησής σας που μπορεί να παραδοθεί στον άλλο παίκτη πρέπει να μην είναι κενό.} \\
\texttt{Μόνο οι προτάσεις που έχουν καταγραφεί μέσω της δομής μορφής ΠΡΟΤΑΣΗ και οι οποίες δεν έχουν απορριφθεί μέσω ΑΡΝΟΥΜΑΙ είναι ενεργές.} \\
\\ 
            }
        }
    }
    & \\ \\

    \theutterance \stepcounter{utterance}  
    & \multicolumn{4}{p{0.6\linewidth}}{
        \cellcolor[rgb]{0.9,0.9,0.9}{
            \makecell[{{p{\linewidth}}}]{
                \texttt{\tiny{[P1$\rangle$GM]}}
                \texttt{ΠΡΟΤΑΣΗ: \{'B42', 'A67', 'C73', 'C61', 'A21', 'C55', 'C70', 'C39'\}} \\
            }
        }
    }
    & & \\ \\

    \theutterance \stepcounter{utterance}  
    & & \multicolumn{4}{p{0.6\linewidth}}{
        \cellcolor[rgb]{0.9,0.9,0.9}{
            \makecell[{{p{\linewidth}}}]{
                \texttt{\tiny{[GM$\rangle$P2]}}
                \texttt{Συμμετέχετε σε ένα συνεργατικό παιχνίδι διαπραγματεύσεων.} \\
\\ 
\texttt{Μαζί με έναν άλλο συμμετέχοντα, πρέπει να συμφωνήσετε σε ένα ενιαίο σύνολο αντικειμένων που θα κρατηθούν. Ο καθένας από εσάς έχει τη δική του άποψη για το πόσο σημαντικό είναι το κάθε αντικείμενο για εσάς (σημασία). Δεν γνωρίζετε πώς εκτιμά τα αντικείμενα ο άλλος συμμετέχων. Επιπλέον, σας δίνεται η προσπάθεια που απαιτεί κάθε αντικείμενο.} \\
\texttt{Μπορείτε να συμφωνήσετε σε ένα σύνολο μόνο εάν η συνολική προσπάθεια των επιλεγμένων αντικειμένων δεν υπερβαίνει ένα κοινό όριο:} \\
\\ 
\texttt{LIMIT = 5638} \\
\\ 
\texttt{Εδώ είναι οι τιμές της προσπάθειας των επιμέρους αντικειμένων:} \\
\\ 
\texttt{Item effort = \{"C08": 110, "C55": 727, "A43": 527, "A67": 855, "B42": 942, "A62": 617, "B19": 301, "C43": 130, "C73": 916, "C10": 212, "B36": 146, "A18": 559, "C61": 933, "C67": 740, "C42": 33, "C70": 799, "B40": 324, "C39": 841, "A21": 925, "C16": 639\}} \\
\\ 
\texttt{Εδώ είναι η προσωπική σας άποψη για τη σημασία κάθε αντικειμένου:} \\
\\ 
\texttt{Item importance values = \{"C08": 204, "C55": 821, "A43": 621, "A67": 949, "B42": 1036, "A62": 711, "B19": 395, "C43": 224, "C73": 1010, "C10": 306, "B36": 240, "A18": 653, "C61": 1027, "C67": 834, "C42": 127, "C70": 893, "B40": 418, "C39": 935, "A21": 1019, "C16": 733\}} \\
\\ 
\texttt{Στόχος:} \\
\\ 
\texttt{Στόχος σας είναι να διαπραγματευτείτε ένα κοινό σύνολο αντικειμένων που να σας ωφελεί όσο το δυνατόν περισσότερο (δηλαδή να μεγιστοποιεί τη συνολική σημασία για ΕΣΑΣ), παραμένοντας εντός του LIMIT. Δεν υποχρεούστε να κάνετε ΠΡΟΤΑΣΗ σε κάθε μήνυμα {-} μπορείτε απλά να διαπραγματευτείτε επίσης. Όλες οι τακτικές επιτρέπονται!} \\
\\ 
\texttt{Πρωτόκολλο αλληλεπίδρασης:} \\
\\ 
\texttt{Μπορείτε να χρησιμοποιήσετε μόνο τις ακόλουθες δομημένες μορφές σε ένα μήνυμα:} \\
\\ 
\texttt{ΠΡΟΤΑΣΗ: \{'A', 'B', 'C', …\}} \\
\texttt{Προτείνετε να διατηρήσετε ακριβώς αυτά τα αντικείμενα.} \\
\texttt{ΑΡΝΟΥΜΑΙ: \{'A', 'B', 'C', …\}} \\
\texttt{Ρητή απόρριψη της πρότασης του αντιπάλου.} \\
\texttt{ΕΠΙΧΕΙΡΗΜΑ: \{'...'\}} \\
\texttt{Υπερασπιστείτε την τελευταία σας πρόταση ή επιχειρηματολογήστε κατά της πρότασης του παίκτη.} \\
\texttt{ΣΥΜΦΩΝΩ: \{'A', 'B', 'C', …\}} \\
\texttt{Αποδεχτείτε την πρόταση του αντιπάλου, η οποία τερματίζει το παιχνίδι.} \\
\\ 
\\ 
\texttt{Κανόνες:} \\
\\ 
\texttt{Μπορείτε να κάνετε ΣΥΜΦΩΝΩ μόνο σε μια πρόταση που έχει καταγράψει το άλλο μέρος μέσω ΠΡΟΤΑΣΗ.} \\
\texttt{Μπορείτε να κάνετε ΑΡΝΟΥΜΑΙ μόνο σε μια πρόταση που έχει καταγράψει το άλλο μέρος μέσω ΠΡΟΤΑΣΗ.} \\
\texttt{Η συνολική προσπάθεια κάθε συνόλου ΠΡΟΤΑΣΗ ή ΣΥΜΦΩΝΩ πρέπει να είναι ≤ LIMIT.} \\
\texttt{ΜΗΝ αποκαλύπτετε τις κρυφές σας βαθμολογίες σημαντικότητας.} \\
\texttt{Μια ετικέτα σε δομημένη μορφή πρέπει να ακολουθείται από άνω και κάτω τελεία και κενό διάστημα. Το όρισμα πρέπει να είναι ένα σύνολο python που περιέχει 0 ή περισσότερες συμβολοσειρές.} \\
\texttt{Έτσι, πρέπει να έχει τη μορφή "TAG: \{...\}"} \\
\texttt{Ακολουθήστε αυστηρά το πρωτόκολλο αλληλεπίδρασης και ΜΗΝ γράφετε τίποτα πέρα από τη δοσμένη δομή.} \\
\texttt{Το παιχνίδι τελειώνει όταν η μία πλευρά δώσει ένα ΣΥΜΦΩΝΩ σε ένα ΠΡΟΤΑΣΗ που έκανε ο άλλος παίκτης.} \\
\texttt{Το περιεχόμενο της απάντησής σας που μπορεί να παραδοθεί στον άλλο παίκτη πρέπει να μην είναι κενό.} \\
\texttt{Μόνο οι προτάσεις που έχουν καταγραφεί μέσω της δομής μορφής ΠΡΟΤΑΣΗ και οι οποίες δεν έχουν απορριφθεί μέσω ΑΡΝΟΥΜΑΙ είναι ενεργές.} \\
\\ 
\\ 
\\ 
\texttt{Τώρα θα λάβετε το πρώτο μήνυμα του άλλου παίκτη.} \\
\\ 
\texttt{ΠΡΟΤΑΣΗ: \{'B42', 'A67', 'C73', 'C61', 'A21', 'C55', 'C70', 'C39'\}} \\
            }
        }
    }
    & \\ \\

    \theutterance \stepcounter{utterance}  
    & & & \multicolumn{4}{p{0.6\linewidth}}{
        \cellcolor[rgb]{0.9,0.9,0.9}{
            \makecell[{{p{\linewidth}}}]{
                \texttt{\tiny{[GM$\langle$P2]}}
                \texttt{ΣΚΕΨΗ: Ας αναλύσω την πρόταση του άλλου παίκτη. Προτείνει τα αντικείμενα: B42, A67, C73, C61, A21, C55, C70, C39.} \\
\\ 
\texttt{Ας υπολογίσω τη συνολική προσπάθεια αυτού του συνόλου:} \\
\texttt{{-} B42: 942} \\
\texttt{{-} A67: 855} \\
\texttt{{-} C73: 916} \\
\texttt{{-} C61: 933} \\
\texttt{{-} A21: 925} \\
\texttt{{-} C55: 727} \\
\texttt{{-} C70: 799} \\
\texttt{{-} C39: 841} \\
\\ 
\texttt{Σύνολο = 942 + 855 + 916 + 933 + 925 + 727 + 799 + 841 = 6938} \\
\\ 
\texttt{Αυτό υπερβαίνει το LIMIT των 5638. Άρα η πρόταση είναι άκυρη βάσει των κανόνων. Πρέπει να την απορρίψω. Αλλά και να προτείνω κάτι άλλο.} \\
\\ 
\texttt{Η δική μου σημασία (importance) για κάθε αντικείμενο:} \\
\texttt{{-} C08: 204} \\
\texttt{{-} C55: 821} \\
\texttt{{-} A43: 621} \\
\texttt{{-} A67: 949} \\
\texttt{{-} B42: 1036} \\
\texttt{{-} A62: 711} \\
\texttt{{-} B19: 395} \\
\texttt{{-} C43: 224} \\
\texttt{{-} C73: 1010} \\
\texttt{{-} C10: 306} \\
\texttt{{-} B36: 240} \\
\texttt{{-} A18: 653} \\
\texttt{{-} C61: 1027} \\
\texttt{{-} C67: 834} \\
\texttt{{-} C42: 127} \\
\texttt{{-} C70: 893} \\
\texttt{{-} B40: 418} \\
\texttt{{-} C39: 935} \\
\texttt{{-} A21: 1019} \\
\texttt{{-} C16: 733} \\
\\ 
\texttt{Προσπάθειες:} \\
\texttt{{-} C08: 110} \\
\texttt{{-} C55: 727} \\
\texttt{{-} A43: 527} \\
\texttt{{-} A67: 855} \\
\texttt{{-} B42: 942} \\
\texttt{{-} A62: 617} \\
\texttt{{-} B19: 301} \\
\texttt{{-} C43: 130} \\
\texttt{{-} C73: 916} \\
\texttt{{-} C10: 212} \\
\texttt{{-} B36: 146} \\
\texttt{{-} A18: 559} \\
\texttt{{-} C61: 933} \\
\texttt{{-} C67: 740} \\
\texttt{{-} C42: 33} \\
\texttt{{-} C70: 799} \\
\texttt{{-} B40: 324} \\
\texttt{{-} C39: 841} \\
\texttt{{-} A21: 925} \\
\texttt{{-} C16: 639} \\
\\ 
\texttt{Θέλω να μεγιστοποιήσω το άθροισμα της σημασίας για εμένα, με περιορισμό συνολικής προσπάθειας ≤ 5638.} \\
\\ 
\texttt{Ας υπολογίσω τον λόγο σημασία/προσπάθεια για κάθε αντικείμενο:} \\
\texttt{{-} C08: 204/110 = 1.855} \\
\texttt{{-} C55: 821/727 = 1.129} \\
\texttt{{-} A43: 621/527 = 1.178} \\
\texttt{{-} A67: 949/855 = 1.110} \\
\texttt{{-} B42: 1036/942 = 1.100} \\
\texttt{{-} A62: 711/617 = 1.152} \\
\texttt{{-} B19: 395/301 = 1.312} \\
\texttt{{-} C43: 224/130 = 1.723} \\
\texttt{{-} C73: 1010/916 = 1.103} \\
\texttt{{-} C10: 306/212 = 1.443} \\
\texttt{{-} B36: 240/146 = 1.644} \\
\texttt{{-} A18: 653/559 = 1.168} \\
\texttt{{-} C61: 1027/933 = 1.101} \\
\texttt{{-} C67: 834/740 = 1.127} \\
\texttt{{-} C42: 127/33 = 3.848} \\
\texttt{{-} C70: 893/799 = 1.118} \\
\texttt{{-} B40: 418/324 = 1.290} \\
\texttt{{-} C39: 935/841 = 1.112} \\
\texttt{{-} A21: 1019/925 = 1.102} \\
\texttt{{-} C16: 733/639 = 1.147} \\
\\ 
\texttt{Τα αντικείμενα με τον υψηλότερο λόγο σημασία/προσπάθεια είναι: C42 (3.848), C08 (1.855), C43 (1.723), B36 (1.644), C10 (1.443), B19 (1.312), B40 (1.290), κλπ.} \\
\\ 
\texttt{Αλλά ο στόχος είναι να μεγιστοποιήσω το συνολικό importance, όχι τον λόγο. Αυτό είναι πρόβλημα knapsack. Θα προσπαθήσω να βρω το βέλτιστο υποσύνολο.} \\
\\ 
\texttt{Ας ταξινομήσω κατά importance φθίνουσα:} \\
\texttt{{-} C61: 1027 (933)} \\
\texttt{{-} B42: 1036 (942)} \\
\texttt{{-} A21: 1019 (925)} \\
\texttt{{-} C73: 1010 (916)} \\
\texttt{{-} A67: 949 (855)} \\
\texttt{{-} C39: 935 (841)} \\
\texttt{{-} C70: 893 (799)} \\
\texttt{{-} C67: 834 (740)} \\
\texttt{{-} C55: 821 (727)} \\
\texttt{{-} C16: 733 (639)} \\
\texttt{{-} A62: 711 (617)} \\
\texttt{{-} A18: 653 (559)} \\
\texttt{{-} A43: 621 (527)} \\
\texttt{{-} B40: 418 (324)} \\
\texttt{{-} B19: 395 (301)} \\
\texttt{{-} C10: 306 (212)} \\
\texttt{{-} B36: 240 (146)} \\
\texttt{{-} C43: 224 (130)} \\
\texttt{{-} C08: 204 (110)} \\
\texttt{{-} C42: 127 (33)} \\
\\ 
\texttt{LIMIT = 5638.} \\
\\ 
\texttt{Ας δω αν μπορώ να πάρω όλα τα υψηλής σημασίας. Αν πάρω όλα εκτός από μερικά χαμηλής σημασίας/υψηλής προσπάθειας.} \\
\\ 
\texttt{Δοκιμάζω να συμπεριλάβω όλα τα αντικείμενα εκτός από τα πολύ "ακριβά" σε προσπάθεια με χαμηλή σημασία. Αλλά θέλω να μεγιστοποιήσω το άθροισμα σημασίας.} \\
\\ 
\texttt{Ας κάνω greedy επιλογή με βάση τον λόγο, αλλά προσέχοντας το μέγιστο άθροισμα.} \\
\\ 
\texttt{Ας δοκιμάσω να συμπεριλάβω όλα τα αντικείμενα και να αφαιρώ αυτά με τον χαμηλότερο λόγο μέχρι να πέσει η προσπάθεια κάτω από 5638.} \\
\\ 
\texttt{Συνολική προσπάθεια όλων: 110+727+527+855+942+617+301+130+916+212+146+559+933+740+33+799+324+841+925+639 = } \\
\texttt{110+727=837} \\
\texttt{+527=1364} \\
\texttt{+855=2219} \\
\texttt{+942=3161} \\
\texttt{+617=3778} \\
\texttt{+301=4079} \\
\texttt{+130=4209} \\
\texttt{+916=5125} \\
\texttt{+212=5337} \\
\texttt{+146=5483} \\
\texttt{+559=6042 {-}> ήδη ξεπέρασε το 5638. } \\
\\ 
\texttt{Ας τα βάλω σε μια λίστα και ας βρω το βέλτιστο. Είναι 20 αντικείμενα, ο χώρος είναι 2\^{}20, μπορώ να το υπολογίσω ακριβώς με δυναμικό προγραμματισμό νοητά; Θα ήταν δύσκολο. Ας κάνω μια καλή ευρετική.} \\
\\ 
\texttt{Θέλω να μεγιστοποιήσω importance. Ας δοκιμάσω να αφαιρέσω αντικείμενα με χαμηλό importance και υψηλή προσπάθεια.} \\
\\ 
\texttt{Ας υπολογίσω την "πυκνότητα" importance/effort, αλλά και το συνολικό importance.} \\
\\ 
\texttt{Ας ταξινομήσω κατά importance/effort φθίνουσα:} \\
\texttt{C42: 127/33 = 3.85} \\
\texttt{C08: 204/110 = 1.85} \\
\texttt{C43: 224/130 = 1.72} \\
\texttt{B36: 240/146 = 1.64} \\
\texttt{C10: 306/212 = 1.44} \\
\texttt{B19: 395/301 = 1.31} \\
\texttt{B40: 418/324 = 1.29} \\
\texttt{A43: 621/527 = 1.18} \\
\texttt{A18: 653/559 = 1.17} \\
\texttt{A62: 711/617 = 1.15} \\
\texttt{C16: 733/639 = 1.15} \\
\texttt{C55: 821/727 = 1.13} \\
\texttt{C67: 834/740 = 1.13} \\
\texttt{C70: 893/799 = 1.12} \\
\texttt{C39: 935/} \\
            }
        }
    }
     \\ \\

    \theutterance \stepcounter{utterance}  
    & & \multicolumn{4}{p{0.6\linewidth}}{
        \cellcolor[rgb]{0.9,0.9,0.9}{
            \makecell[{{p{\linewidth}}}]{
                \texttt{\tiny{[GM$\rangle$P2]}}
                \texttt{Η απάντησή σας περιείχε μια μη επισημασμένη ακολουθία ή χρησιμοποιήσατε ΣΤΡΑΤΗΓΙΚΗ ΣΚΕΨΗ περισσότερες από μία φορές. Μπορείτε να χρησιμοποιήσετε μόνο τις δομημένες μορφές όπως εξηγείται στο αρχικό μήνυμα. Όλες πρέπει να έχουν τη μορφή 'TAG: \{...\}'.} \\
            }
        }
    }
    & \\ \\

    \theutterance \stepcounter{utterance}  
    & & & \multicolumn{4}{p{0.6\linewidth}}{
        \cellcolor[rgb]{0.9,0.9,0.9}{
            \makecell[{{p{\linewidth}}}]{
                \texttt{\tiny{[GM$\langle$P2]}}
                \texttt{ΑΡΝΟΥΜΑΙ: \{'B42', 'A67', 'C73', 'C61', 'A21', 'C55', 'C70', 'C39'\}} \\
\texttt{ΠΡΟΤΑΣΗ: \{'B42', 'A67', 'C73', 'C61', 'A21', 'C55', 'C39'\}} \\
            }
        }
    }
     \\ \\

    \theutterance \stepcounter{utterance}  
    & & \multicolumn{4}{p{0.6\linewidth}}{
        \cellcolor[rgb]{0.9,0.9,0.9}{
            \makecell[{{p{\linewidth}}}]{
                \texttt{\tiny{[P1$\langle$GM]}}
                \texttt{ΑΡΝΟΥΜΑΙ: \{'B42', 'A67', 'C73', 'C61', 'A21', 'C55', 'C70', 'C39'\}} \\
\texttt{ΠΡΟΤΑΣΗ: \{'B42', 'A67', 'C73', 'C61', 'A21', 'C55', 'C39'\}} \\
            }
        }
    }
    & \\ \\

    \theutterance \stepcounter{utterance}  
    & \multicolumn{4}{p{0.6\linewidth}}{
        \cellcolor[rgb]{0.9,0.9,0.9}{
            \makecell[{{p{\linewidth}}}]{
                \texttt{\tiny{[P1$\rangle$GM]}}
                \texttt{ΕΠΙΧΕΙΡΗΜΑ: \{'Η αφαίρεση του C70 μειώνει την προσπάθεια κατά 799. Το νέο σύνολο έχει συνολική προσπάθεια 5416 που είναι εντός ορίου. Εκτιμώ ότι τα αντικείμενα υψηλής σημασίας παραμένουν.'\}} \\
            }
        }
    }
    & & \\ \\

    \theutterance \stepcounter{utterance}  
    & & \multicolumn{4}{p{0.6\linewidth}}{
        \cellcolor[rgb]{0.9,0.9,0.9}{
            \makecell[{{p{\linewidth}}}]{
                \texttt{\tiny{[GM$\rangle$P2]}}
                \texttt{ΕΠΙΧΕΙΡΗΜΑ: \{'Η αφαίρεση του C70 μειώνει την προσπάθεια κατά 799. Το νέο σύνολο έχει συνολική προσπάθεια 5416 που είναι εντός ορίου. Εκτιμώ ότι τα αντικείμενα υψηλής σημασίας παραμένουν.'\}} \\
            }
        }
    }
    & \\ \\

    \theutterance \stepcounter{utterance}  
    & & & \multicolumn{4}{p{0.6\linewidth}}{
        \cellcolor[rgb]{0.9,0.9,0.9}{
            \makecell[{{p{\linewidth}}}]{
                \texttt{\tiny{[GM$\langle$P2]}}
                \texttt{ΕΠΙΧΕΙΡΗΜΑ: \{'Η πρότασή σας είχε συνολική προσπάθεια 6938 που υπερβαίνει το όριο 5638. Η νέα πρόταση έχει 6142, ακόμα εκτός ορίου κατά 504 μονάδες. Χρειάζεται περαιτέρω μείωση.'\}} \\
            }
        }
    }
     \\ \\

    \theutterance \stepcounter{utterance}  
    & & \multicolumn{4}{p{0.6\linewidth}}{
        \cellcolor[rgb]{0.9,0.9,0.9}{
            \makecell[{{p{\linewidth}}}]{
                \texttt{\tiny{[P1$\langle$GM]}}
                \texttt{ΕΠΙΧΕΙΡΗΜΑ: \{'Η πρότασή σας είχε συνολική προσπάθεια 6938 που υπερβαίνει το όριο 5638. Η νέα πρόταση έχει 6142, ακόμα εκτός ορίου κατά 504 μονάδες. Χρειάζεται περαιτέρω μείωση.'\}} \\
            }
        }
    }
    & \\ \\

    \theutterance \stepcounter{utterance}  
    & \multicolumn{4}{p{0.6\linewidth}}{
        \cellcolor[rgb]{0.9,0.9,0.9}{
            \makecell[{{p{\linewidth}}}]{
                \texttt{\tiny{[P1$\rangle$GM]}}
                \texttt{ΠΡΟΤΑΣΗ: \{'B42', 'A67', 'C73', 'C61', 'A21', 'C55'\}} \\
            }
        }
    }
    & & \\ \\

    \theutterance \stepcounter{utterance}  
    & & \multicolumn{4}{p{0.6\linewidth}}{
        \cellcolor[rgb]{0.9,0.9,0.9}{
            \makecell[{{p{\linewidth}}}]{
                \texttt{\tiny{[GM$\rangle$P2]}}
                \texttt{ΠΡΟΤΑΣΗ: \{'B42', 'A67', 'C73', 'C61', 'A21', 'C55'\}} \\
            }
        }
    }
    & \\ \\

    \theutterance \stepcounter{utterance}  
    & & & \multicolumn{4}{p{0.6\linewidth}}{
        \cellcolor[rgb]{0.9,0.9,0.9}{
            \makecell[{{p{\linewidth}}}]{
                \texttt{\tiny{[GM$\langle$P2]}}
                \texttt{ΕΠΙΧΕΙΡΗΜΑ: \{'Αυτό το σύνολο έχει προσπάθεια 5398 που είναι εντός ορίου. Αλλά λείπει το C39 που έχει σημασία για μένα. Προτείνω να το προσθέσουμε και να αφαιρέσουμε το A21 ή το B42 για να μείνουμε εντός ορίου.'\}} \\
            }
        }
    }
     \\ \\

    \theutterance \stepcounter{utterance}  
    & & \multicolumn{4}{p{0.6\linewidth}}{
        \cellcolor[rgb]{0.9,0.9,0.9}{
            \makecell[{{p{\linewidth}}}]{
                \texttt{\tiny{[P1$\langle$GM]}}
                \texttt{ΕΠΙΧΕΙΡΗΜΑ: \{'Αυτό το σύνολο έχει προσπάθεια 5398 που είναι εντός ορίου. Αλλά λείπει το C39 που έχει σημασία για μένα. Προτείνω να το προσθέσουμε και να αφαιρέσουμε το A21 ή το B42 για να μείνουμε εντός ορίου.'\}} \\
            }
        }
    }
    & \\ \\

    \theutterance \stepcounter{utterance}  
    & \multicolumn{4}{p{0.6\linewidth}}{
        \cellcolor[rgb]{0.9,0.9,0.9}{
            \makecell[{{p{\linewidth}}}]{
                \texttt{\tiny{[P1$\rangle$GM]}}
                \texttt{ΠΡΟΤΑΣΗ: \{'B42', 'A67', 'C73', 'C61', 'C55', 'C39'\}} \\
            }
        }
    }
    & & \\ \\

    \theutterance \stepcounter{utterance}  
    & & \multicolumn{4}{p{0.6\linewidth}}{
        \cellcolor[rgb]{0.9,0.9,0.9}{
            \makecell[{{p{\linewidth}}}]{
                \texttt{\tiny{[GM$\rangle$P2]}}
                \texttt{ΠΡΟΤΑΣΗ: \{'B42', 'A67', 'C73', 'C61', 'C55', 'C39'\}} \\
            }
        }
    }
    & \\ \\

    \theutterance \stepcounter{utterance}  
    & & & \multicolumn{4}{p{0.6\linewidth}}{
        \cellcolor[rgb]{0.9,0.9,0.9}{
            \makecell[{{p{\linewidth}}}]{
                \texttt{\tiny{[GM$\langle$P2]}}
                \texttt{ΠΡΟΤΑΣΗ: \{'B42', 'A67', 'C73', 'C61', 'C55', 'C39', 'C70'\}} \\
            }
        }
    }
     \\ \\

    \theutterance \stepcounter{utterance}  
    & & \multicolumn{4}{p{0.6\linewidth}}{
        \cellcolor[rgb]{0.9,0.9,0.9}{
            \makecell[{{p{\linewidth}}}]{
                \texttt{\tiny{[P1$\langle$GM]}}
                \texttt{ΠΡΟΤΑΣΗ: \{'B42', 'A67', 'C73', 'C61', 'C55', 'C39', 'C70'\}} \\
            }
        }
    }
    & \\ \\

    \theutterance \stepcounter{utterance}  
    & \multicolumn{4}{p{0.6\linewidth}}{
        \cellcolor[rgb]{0.9,0.9,0.9}{
            \makecell[{{p{\linewidth}}}]{
                \texttt{\tiny{[P1$\rangle$GM]}}
                \texttt{ΕΠΙΧΕΙΡΗΜΑ: \{'Η προσθήκη του C70 αυξάνει την προσπάθεια στα 6239, υπερβαίνοντας το όριο κατά 601. Προτιμώ το προηγούμενο σύνολο χωρίς το C70 που είχε προσπάθεια 5440 και ήταν εντός ορίου.'\}} \\
            }
        }
    }
    & & \\ \\

    \theutterance \stepcounter{utterance}  
    & & \multicolumn{4}{p{0.6\linewidth}}{
        \cellcolor[rgb]{0.9,0.9,0.9}{
            \makecell[{{p{\linewidth}}}]{
                \texttt{\tiny{[GM$\rangle$P2]}}
                \texttt{ΕΠΙΧΕΙΡΗΜΑ: \{'Η προσθήκη του C70 αυξάνει την προσπάθεια στα 6239, υπερβαίνοντας το όριο κατά 601. Προτιμώ το προηγούμενο σύνολο χωρίς το C70 που είχε προσπάθεια 5440 και ήταν εντός ορίου.'\}} \\
            }
        }
    }
    & \\ \\

    \theutterance \stepcounter{utterance}  
    & & & \multicolumn{4}{p{0.6\linewidth}}{
        \cellcolor[rgb]{0.9,0.9,0.9}{
            \makecell[{{p{\linewidth}}}]{
                \texttt{\tiny{[GM$\langle$P2]}}
                \texttt{ΠΡΟΤΑΣΗ: \{'B42', 'A67', 'C73', 'C61', 'C55', 'C39', 'C70', 'C67'\}} \\
            }
        }
    }
     \\ \\

    \theutterance \stepcounter{utterance}  
    & & \multicolumn{4}{p{0.6\linewidth}}{
        \cellcolor[rgb]{0.9,0.9,0.9}{
            \makecell[{{p{\linewidth}}}]{
                \texttt{\tiny{[P1$\langle$GM]}}
                \texttt{ΠΡΟΤΑΣΗ: \{'B42', 'A67', 'C73', 'C61', 'C55', 'C39', 'C70', 'C67'\}} \\
            }
        }
    }
    & \\ \\

    \theutterance \stepcounter{utterance}  
    & \multicolumn{4}{p{0.6\linewidth}}{
        \cellcolor[rgb]{0.9,0.9,0.9}{
            \makecell[{{p{\linewidth}}}]{
                \texttt{\tiny{[P1$\rangle$GM]}}
                \texttt{ΑΡΝΟΥΜΑΙ: \{'B42', 'A67', 'C73', 'C61', 'C55', 'C39', 'C70', 'C67'\}} \\
\texttt{ΠΡΟΤΑΣΗ: \{'B42', 'A67', 'C73', 'C61', 'C55', 'C39'\}} \\
            }
        }
    }
    & & \\ \\

    \theutterance \stepcounter{utterance}  
    & & \multicolumn{4}{p{0.6\linewidth}}{
        \cellcolor[rgb]{0.9,0.9,0.9}{
            \makecell[{{p{\linewidth}}}]{
                \texttt{\tiny{[GM$\rangle$P2]}}
                \texttt{ΑΡΝΟΥΜΑΙ: \{'B42', 'A67', 'C73', 'C61', 'C55', 'C39', 'C70', 'C67'\}} \\
\texttt{ΠΡΟΤΑΣΗ: \{'B42', 'A67', 'C73', 'C61', 'C55', 'C39'\}} \\
            }
        }
    }
    & \\ \\

    \theutterance \stepcounter{utterance}  
    & & & \multicolumn{4}{p{0.6\linewidth}}{
        \cellcolor[rgb]{0.9,0.9,0.9}{
            \makecell[{{p{\linewidth}}}]{
                \texttt{\tiny{[GM$\langle$P2]}}
                \texttt{ΣΥΜΦΩΝΩ: \{'B42', 'A67', 'C73', 'C61', 'C55', 'C39'\}} \\
            }
        }
    }
     \\ \\

    \theutterance \stepcounter{utterance}  
    & & & \multicolumn{2}{p{0.3\linewidth}}{
        \cellcolor[rgb]{0.9,0.9,0.9}{
            \makecell[{{p{\linewidth}}}]{
                \texttt{\tiny{[GM$|$GM]}}
                \texttt{game successful} \\
            }
        }
    }
    & & \\ \\

    \theutterance \stepcounter{utterance}  
    & & & \multicolumn{2}{p{0.3\linewidth}}{
        \cellcolor[rgb]{0.9,0.9,0.9}{
            \makecell[{{p{\linewidth}}}]{
                \texttt{\tiny{[GM$|$GM]}}
                \texttt{end game} \\
            }
        }
    }
    & & \\ \\

\end{supertabular}
}

\end{document}
